\documentclass[10pt]{article}
\usepackage[margin=20mm]{geometry}
\usepackage{amsmath,amssymb,mathtools,fontspec,unicode-math,xurl,hyperref}
\IfFontExistsTF{DejaVu Serif}{\setmainfont{DejaVu Serif}}{\setmainfont{Latin Modern Roman}}
\setmathfont{latinmodern-math.otf}
\setlength{\parindent}{0pt}\setlength{\parskip}{6pt}\sloppy
\title{Chaos, spectral measures and the meaning of dimension}
\author{AI-assisted technical review}
\date{}
\begin{document}
\maketitle
Abstract: The principal identification of the circular Gaussian multiplicative chaos with a randomly weighted spectral measure was checked through its normalization, change of measure, diffusion limit, entrance law, and conditional limiting argument. No substantive defect in that identification was established in the inspected scope. Corollary 2.4 needs a precise distinction between the Hausdorff dimension of a measure and that of its topological support. The latter is the full circle for the stated Verblunsky law, including the critical case. This is a Minor correction to the formulation of the dimensional consequence, not a refutation of the measure correspondence. Actual page inspection also confirms a reversed probability-event sign in the entrance-law proof; the preceding correct tail split supplies its local repair.\par
\section*{Bound source and scope}
This report concerns Reda Chhaibi and Joseph Najnudel, On the circle, Gaussian multiplicative chaos and Beta ensembles match exactly, DOI 10.4171/JEMS/1700 [paper]. The bound source is the 77-page publisher PDF, serial-article file 52155, with SHA256\par
`3cdc097f311e31d65ef6ad959800cc52cc15`
`e1543fd165e6d5ae9f365337027d`\par
The hash is recorded without abbreviation in the accompanying provenance and inventory. Its first page gives J. Eur. Math. Soc. 28, 1–77 (2026) and copyright 2025; the publisher records publication on 7 November 2025 and acceptance on 14 February 2025. The first page records receipt on 1 September 2019 and revision on 18 January 2024. These facts do not change the frozen identifier M25-0324 or its assigned 2025 collection. No substitution of an arXiv revision or an online-first PDF has been made.\par
All 77 source pages, including Appendix A and the references, were read as extracted text. In the same-chat continuation, all 25 required source evidence images were actually displayed and inspected. This confirms the support notation and resolves the previously unconfirmed probability-event sign on page 68. Earlier `EMFILE` failures and false viewing records are retained. The public publisher text remains secondary corroboration; the exact bound PDF supplies the actual visual evidence.\par
\section*{Objects, parameters, and dependencies}
Write \(q=2/\beta=\gamma^2\). The independent Verblunsky coefficients have rotation-invariant laws with \[|\alpha_j|^2\sim{\rm Beta}(1,(j+1)/q).\] The limiting probability measure \(\mu^\beta\) uses this infinite coefficient sequence. The finite approximation involves random spectral weights after reversing the finite coefficient order; it is not the unweighted empirical eigenvalue measure. Equality of the finite spectral point law does not identify those weights with the usual independent Dirichlet weights.\par
The real Gaussian field has covariance \(-2\log|1-w\bar z|\). Thus the radial subcritical regularization has factor \((1-r^2)^{\gamma^2}\), with normalized angular measure. For \(\beta>2\), the total-mass factor is \[C_0=\prod_{j\ge0}(1-|\alpha_j|^2)^{-1}
                       \left(1-\frac{q}{j+1}\right).\] It is a functional of the same coefficients as \(\mu^\beta\); independence of mass and normalized measure is not asserted. At \(\beta=2\) the first deterministic factor vanishes, and the theorem uses the separately renormalized expression with leading factor \(2\). The critical chaos is defined through the subcritical limiting normalization. Applying the subcritical product without that change would be incorrect. The supercritical comparison in Section 2.5 is explicitly prospective; the main identification does not prove equivalence with either freezing or KPZ-dual constructions.\par
The proof’s substantial bridge is not merely convergence of the analytic polynomial inside the disk. It is the weighted \(L^1\) estimate of Lemma 4.1 and the conditional expectation limit of Lemma 4.2. These are supplied only in Section 7 after the estimates of Section 5 and the diffusion analysis of Section 6. The distinction is essential because radial and polynomial-degree limits cannot be exchanged on the basis of compact interior convergence alone.\par
\section*{A dimensional formulation requiring correction}
Problem 1 (Minor: Corollary 2.4, page 11).  With the conventional meaning of topological support, the displayed equalities \(\dim_H\operatorname{supp}(\mu^\beta)=1-2/\beta\) for \(\beta>2\), and \(\dim_H\operatorname{supp}(\mu^2)=0\), are false. The intended measure dimension must be defined and used explicitly. The analogous GMC displays require the same distinction.\par
\subsection*{A counterexample within the stated law}
Take \(\beta=4\), so all of the subcritical assumptions hold. For any \(0<\varepsilon<1\), the exact beta tail is \[\mathbb P(|\alpha_j|>\varepsilon)
   =(1-\varepsilon^2)^{2(j+1)}.\] The sum over \(j\) is finite. Borel–Cantelli and a countable sequence of \(\varepsilon\)’s give \(\alpha_j\to0\) almost surely. This argument works for every \(\beta>0\), not only for the illustrative value \(4\).\par
Here is a direct spectral reason why that implies full topological support. In the CMV representation used on pages 19–21, each \(2\)-by-\(2\) block differs from its zero-coefficient block by entries tending to zero: the diagonal entries are \(\alpha_j\) and its conjugate, up to signs, and the off-diagonal difference is \(\sqrt{1-|\alpha_j|^2}-1\). Consequently the difference between the corresponding infinite CMV unitary and the free CMV unitary is compact; indeed the block-diagonal differences are norm limits of their finite truncations.\par
For completeness, the free matrix has every \(\zeta\in\partial\mathbb D\) in its essential spectrum. With zero-based basis indexing its odd tail satisfies \(C_{\rm free}e_{2k+1}=e_{2k+3}\). For \(J\) beyond the boundary and \(N\ge1\), put \[v_{J,N}=\frac1{\sqrt N}
       \sum_{k=J}^{J+N-1}\zeta^{-k}e_{2k+1}.\] The two endpoint errors give exactly \[\|(C_{\rm free}-\zeta)v_{J,N}\|^2=\frac2N.\] Taking the blocks farther out makes the perturbation error arbitrarily small, since \(\alpha_j\to0\). These are Weyl sequences for every \(\zeta\). The spectrum is therefore the whole unit circle. In the cyclic spectral representation of the CMV matrix, the spectrum is the closed support of its spectral measure. Hence, almost surely, \[\operatorname{supp}_{\rm top}\mu^\beta=\partial\mathbb D,\qquad
 \dim_H\operatorname{supp}_{\rm top}\mu^\beta=1 .\] For \(\beta=4\) the displayed value in Corollary 2.4 is \(1/2\). For \(\beta=2\) it is \(0\). The discrepancy is thus an actual almost-sure counterexample under the prescribed distributions, not a pathological choice of coefficients outside the assumptions.\par
\subsection*{Impact and repair boundary}
A singular measure can have a dense carrier of small Hausdorff dimension and nevertheless have full closed support. An appropriate quantity to state here is, for example, \[\dim_H\mu=\inf\{\dim_H E:E\text{ Borel},\ \mu(E)=1\}.\] The intended assertions about this quantity, and any stronger exact-dimensional interpretation, should be tied explicitly to the spectral theorem cited by the authors. This report does not independently reprove Simon’s Theorem 12.7.7 or the extension of its antecedent hypotheses mentioned in the paragraph following the corollary.\par
The finding is classified as Minor because it concerns specification of the dimensional object in a consequence. It neither invalidates the main random-measure identity nor establishes a false value of the intended measure dimension. The phrase and notation should nevertheless be changed: no nonstandard definition of \(\operatorname{supp}\) resolving the distinction was located in the full text.\par
\section*{Checks of the main proof}
\subsection*{Analytic martingales and change of measure}
The logarithm of the reversed polynomial is based at zero; its nonvanishing inside the disk permits that branch. Rotational independence gives the martingale increments mean zero. Proposition 3.6 uses the exponential cubic-sum hypothesis and a conditional bounded-differences estimate after conditioning on moduli. The trace expansion in Lemma 3.5 provides a cubic remainder with an absolute bound. The conditional quadratic dependence in the phase variables is summable under the stated fourth-power bound, which is itself covered by the hypothesis. The Hilbert-space argument uses moment bounds and holomorphy to pass from Fourier-coefficient convergence to convergence on compact disks. Identification of the Gaussian coefficients uses finite-dimensional trace limits, with the cited result or Appendix A as an external symmetric-function input.\par
For fixed \(r<1\), the density in (4.7) is a genuine positive martingale, bounded in \(L^p\) by the already established polynomial moment bounds. The conditional angular average of the Poisson kernel is one. This verifies the change-of-measure step and explains why the moduli retain their independent beta laws under \(\mathbb Q_{r,\theta}\). The subsequent limit concerns this \(r\)-dependent law; it is not an unjustified application of a theorem under a single fixed probability measure.\par
\subsection*{Moment and tail estimates}
The angular series in Lemma 5.2 was checked against the constant and quadratic terms. Under the tilted law, the conditional drift and variance of the modulus recursion give the coefficients in (5.6)–(5.8). The fourth-moment cancellation uses the alternating binomial coefficients \(1,-4,6,-4,1\). The early finitely many beta variables, whose inverse moments may diverge, are treated by truncation rather than by extending a large-index moment bound to every index.\par
In Proposition 5.5, the logarithmic error splits into a centered martingale and a nonpositive remainder. The concavity inequality for \(\log(1-aq)\) supplies that sign. The uniform exponential bound uses positive moments of that error, not nonexistent negative moments. The tail proof in Proposition 5.4 is naturally an \(A\ge1\) argument. For \(0<A<1\), the additional initial sum is bounded deterministically by a harmonic sum of size \(O(1+\log(1/A))\), which is \(O(1/A)\); this supplies the harmless small-\(A\) extension. No core failure is inferred from the compressed treatment of that range.\par
\subsection*{Diffusion, entrance, and uniqueness}
The mesh is \(\log(1/r^2)\), with \(1-r^2\) asymptotically equal to it. On intervals \([\varepsilon,A]\), the coefficients have the required uniform convergence and the fourth-moment estimates imply tightness. The coefficients only need bounded extensions outside the invariant state interval \([0,1]\); the displayed quadratic drift should not be mistaken for a bounded function on all of \(\mathbb R\) without such an extension. The earlier tightness estimate near zero supplies the additional entrance restriction \(1-X_t=O(t^{1-\varepsilon_0})\). The SDE on positive times alone does not select that entrance.\par
Put \(Y_t=1-X_t\) and \(L_t=\log(Y_t/t)\). Direct Itô differentiation gives \[dL_t=\left\{\frac{e^{-L_t}+2/\beta-1}{t}
                      -1-\frac4\beta e^{L_t}\right\}dt
       +\sqrt{\frac{4(1-te^{L_t})}{\beta t}}\,dB_t .\] The derivative of the drift with respect to \(L\) is \(-e^{-L}/t-4e^L/\beta<0\). This independently checks the monotonicity used in Lemma 6.4. Together with noncrossing of synchronously coupled solutions on positive time intervals, it yields the required contraction estimate. The Volterra transformation and \(s=t e^{-\beta w}\) give the inverse-gamma entrance parameter \(\beta/2-1>0\); this is another point at which subcriticality is essential. No claim of applying that entrance proof directly at \(\beta=2\) is justified.\par
Problem 2 (Minor: the first probability display on page 68).  Actual inspection confirms that the first probability on the right uses the lower-tail event \(I_t\leq\alpha\), where \[I_t=\int_0^t\frac{4(1-X_u)}{\beta u}\,du.\] The last display on page 67 instead has the required upper-tail event \(I_t\geq\alpha\).\par
Under precisely the stated near-zero assumption, choose \(\varepsilon_0=1/2\). Then \(0\leq1-X_u\leq C(\omega)\sqrt u\) for small \(u\), with \(C(\omega)<\infty\) almost surely. Direct integration gives \[0\leq I_t\leq\frac{8C(\omega)}{\beta}\sqrt t\longrightarrow0.\] For every fixed \(\alpha>0\), bounded convergence therefore yields \(\mathbb P(I_t\leq\alpha)\to1\), whereas \(\mathbb P(I_t\geq\alpha)\to0\). Thus the lower-tail term printed on page 68 cannot be dropped in the subsequent limit. Replacing its sign by the upper-tail sign already present on page 67 restores the inference; the remaining Brownian supremum tail tends to zero as \(\alpha\downarrow0\). This is a confirmed local transcription correction, not a counterexample to the entrance law. The exact integral was independently checked symbolically, and the analytic probability limits above require no simulation or author code.\par
\subsection*{Conditional limit and mass}
The last section first removes the logarithmic error using uniform exponential integrability, then truncates the exponential of the remaining variable. The sign-event decomposition on page 72 permits conditioning events to depend on \(r\), which is explicitly allowed in Proposition 6.1. Thus weak convergence for those conditioned processes, with small- and large-time tails controlled separately, proves the needed \(L^1\) conditional convergence. Ordinary unconditioned weak convergence would not suffice.\par
For a finite product with \(U_n=1-|\alpha_{n-1}|^2\), \(U_n\sim{\rm Beta}(n/q,1)\), an independent beta integral gives \[\mathbb E[C_N^z]
   =\prod_{n=1}^N\frac{(1-q/n)^z}{1-qz/n}.\] Its gamma-product limit is \[\frac{\Gamma(1-qz)}{\Gamma(1-q)^z},\] in the stated safe range of Mellin parameters. This agrees with the claimed inverse-exponential total mass and does not assume independence between the mass and normalized measure. The critical prefactor follows from \(\Gamma(1-\gamma^2)^{-1}/(1-\gamma)\to2\). An exact symbolic check was executed for the Itô drift and its derivative; a rational finite-chain check confirmed the two-endpoint Weyl-sequence norm. These are reviewer-created checks, not author code, and the analytic support argument is not inferred from numerical samples.\par
\section*{Coverage and limitations}
The introduction and corollaries (pages 1–15), analytic and trace arguments (16–27), conditional proof framework (27–33), estimates (34–47), diffusion and entrance proofs (47–70), final limiting argument (70–74), and Appendix A and references (74–77) were read. No independent reproof of the full OPUC correspondence, all cited critical-chaos results, Dufresne’s identity, or the Macdonald scalar-product theorem is claimed. The detailed ledger identifies these dependencies.\par
A bounded title/author search and the publisher record did not locate a verified paper-specific supplementary file, code repository, or erratum. Absence of retrieved code is not a mathematical defect, and no author code was executed. No supplementary material is declared absent merely because it was not located.\par
The present scientific assessment is Minor, with no confirmed Major finding in the inspected scope. It remains an initial technical review, not independent acceptance. The machine-readable `scientific\_decision` remains `pending`, human review fields remain empty. The 25 source evidence pages have now been actually inspected. The prior report and all failed or false records are preserved. This revised TeX must be freshly compiled and every resulting PDF page actually inspected before completion; the continuation provenance binds the current source, TeX, compiler, dependencies, PDF and observed images.\par
9 R. Chhaibi and J. Najnudel, On the circle, Gaussian multiplicative chaos and Beta ensembles match exactly, J. Eur. Math. Soc. 28 (2026), 1–77, DOI 10.4171/JEMS/1700. Bound publisher PDF: https://ems.press/content/serial-article-files/52155?nt=1. EMS Press, article and version record, https://ems.press/journals/jems/articles/14299254, accessed during this review.\par
\end{document}
